\section{THE BASIC GAME RULES}

\subsection{INTRODUCTION}

The Basic Game is the easiest to learn and least complicated version of SHENANDOAH. It is ideal for players who prefer relatively simple and straightforward games, and also serves as an introduction to the game's mechanics for players who prefer more complex games and will eventually wish to try the Advanced and Optional versions of the Game.

Unlike many other game simulations that more experienced players may have tried in the past where movement and battle are broken down into separate steps; in SHENANDOAH movement and battle are often intermingled, with individual units being capable of moving, fighting, moving again, fighting again, etc, during the course of the Turn.

Movement in SHENANDOAH is simulated by moving the unit counters from hex to hex across the gridded mapboard. Each hex entered costs the unit being moved a portion of its Movement Allowance Factor, as does regrouping the units (Stacking and Unstacking them), changing their Formation, and engaging in Combat. A unit's Movement Allowance Factor is a numerical representation of its movement abilities, and the cost for various movements simulates the relative difficulties of these actions.

Combat in SHENANDOAH is simulated by comparing the relative strengths of the forces involved, their Formation, and the strength of their position. The often fluid nature of the Combat in the actual campaign is reflected in the Combat simulation mechanics of the Game.

\subsection{OUTLINE OF PLAY}

Each scenario in SHENANDOAH consists of a certain number of Turns. When that number of Turns is over, the game is over. Each Turn is divided into two parts; the Union Player's portion of the Turn, and the Cofederate Player's portion of the Turn, not necessarily in that order. The player who is to move first during each Turn is determined in the individual scenarios. The player who is currently in his portion, or Phase, of the the Turn is referred to as the Phasing Player, and the other player is referred to as the Non-Phasing Player. Below is a step-by-step outline of the sequence of events during a turn.

\begin{enumerate}
  \item The player whose side moves first as determined in the individual scenario is defined now as the Phasing player, and his opponent is defined as the Non-Phasing Player. The Phasing Player may now move all of his units in any direction up to the limit of their Movement Allowance Factors, within the limits and restrictions of the rules of movement. The Phasing Player should note on a scrap of scratch paper the number of Movement Allowance Factors expended by units or Stacks of units that have entered enemy Zones of Control. This information is \textit{very important} as these units in enemy Zones of Control can use any remaining Movement Allowance Factors for Combat and/or further movement.

  \item The Non-Phasing Player is now allowed to exercise his movement options, either to fall back to avoid Combat, or to stay in place and accept it.

  \item Any Combat situations are now resolved. Combat costs the Phasing Player Movement Allowance Factors, and each Combat situation continues until the Phasing Player's units use up their Movement Allowance Factors, or until the Combat is resolved to the point where the Phasing Player's units are no longer in an enemy Zone of Control.

  \item If the Phasing Player still has units with unexpended Movement Allowance Factors, he may use them for further movement and/or combat. He is allowed again to enter enemy Zones of Control, if desired, again noting on his scratch paper the current number of Movement Allowance Factors expended.

  \item Steps 2-4 are repeated until the Phasing Player has no Movement Allowance Factors left, or has no further movement he wishes to make.

  \item The player who moved first has now completed his portion of the Turn. The player moving second is now defeind to be the Phasing Player, and his opponent becomes the Non-Phasing Player. Steps 1-5 are now repeated for this player's portion of the Turn.

  \item The players now record on their Point Indicators their new point totals.

  \item The Turn Marker is advanced on the Game Turn track, signaling the completion of one Turn, and the start of another.
\end{enumerate}

\subsection{ZONES OF CONTROL}

\begin{enumerate}
  \item Infantry, Cavalry, Partisan, Artillery, and Horse Artillery units all have a Zone of Control. No other types of units in the game have Zones of Control.

  \item The six hexes that immediately surround a hex and are adjacent to it constitute the Zone of Control of any unit or units in that hex.

  \item A Zone of Control does not extend into a hex across a ridgeline.

  \item All units must stop as soon as they are moved into the Zone of Control of an enemy unit - movement can be continued only after the enemy unit has moved back (either voluntarily or involuntarily, as a result of combat) and no enemy Zones of Control cover the hex.

  \item Zones of Control do not effect combat after an initial round of combat has occurred. Enemy units in each other's Zones of Control do not necessarily have to have combat.

  \item Units may retreat into enemy Zones of Control, but only if no alternate retreat routes are available.

  \item A unit or Stack of units cannot move into an enemy Zone of Control if it lacks enough unexpended Movement Allowance Factors to engage in Combat after entering.
\end{enumerate}

(NOTE: When any of the Phasing Player's units enter an enemy zone as a result of voluntary movement they are obligated to initiate one round of combat - as explained more fully later in the rules.)

\subsection{ENTERING AND EXITING THE MAPBOARD}

\begin{enumerate}
  \item Each scenario lists the hexes where the various reinforcement units may enter the Mapboard. They may enter on those hexes, or within four hexes of the designated hex, along the same edge of the Mapboard. The player does not have to bring these units into the game; it is his choice to do so or not. The decision as to whether the units will be brought on must be made by the player ast the start of the game (exceptions: In the "long" scenarios \#6 and \#16, the player must make his decision seven Turns before the troops will arrive.); this decision to bring the troops on is indicated by placing the units to be brought on the proper Turn of their arrival on the TURN RECORD CHART.

  \item Not all reinforcement units need to be brought on; if the player wishes he may commit only part of a force, as indicated in the scenarios. Once the player decides to bring the units on, they must be placed on the Mapboard when their turn arrives, and remain for at least that Turn.

  \item Units may exit the Mapboard at certain designated hexes (see each individual scenario). Units that exit on these designated hexes are allowed to re-enter the Mapboard on a later Turn. This re-entry must be at the hex where they exited, or within four hexes of it along the same edge of the Mapboard. The decision to bring units back (again placing them in the proper box of the TURN RECORD CHART) must be made a number of Turns prior to their re-entry; that number of Turns being equal to the number of Turns since ithe units exited the Mapboard. For instance, if on Turn Four a player decides to re-enter units that exited on Turn Two, their re-entry Turn would be Turn Six.
\end{enumerate}